\begin{proof}[Proof of \cref{obj:lem:parametric-null-floor}]
\begin{enumerate}
\item Fix the parameters in the statement. For every \(\tau\in\mathbb R\), define the observed law \(P_{\mathrm{par},\tau}\) by
\[
P_{\mathrm{par},\tau}(dx,\{(a,y)\})
=
\lambda(dx)
\begin{cases}
1/4,&a=0,\ y\in\{0,1\},\\
\frac12\{1-\ell(\tau)\},&a=1,\ y=0,\\
\frac12\ell(\tau),&a=1,\ y=1,
\end{cases}
\qquad
x\in[0,1],\quad a,y\in\{0,1\},
\]
where \(\ell\) is the logistic map,
\[
\ell(\tau)=\frac{1}{1+\exp(-\tau)}.
\]
The four conditional cell probabilities are constant in \(x\), strictly positive, and sum to one. Thus this construction has continuous interior cells and uniform covariate marginal \(\lambda\), which has full support.

Summing the treated cells gives treatment probability \(1/2\); dividing each success cell by its arm probability gives control risk \(1/2\) and treated risk \(\ell(\tau)\). Moreover,
\[
\frac{\ell(\tau)}{1-\ell(\tau)}=\exp(\tau),
\qquad
\operatorname{logit}\{\ell(\tau)\}=\tau,
\qquad
\operatorname{logit}(1/2)=0.
\]
Consequently, for every \(x\in[0,1]\),
\[
\begin{aligned}
e(P_{\mathrm{par},\tau},x)&=\frac12,\\
g(P_{\mathrm{par},\tau},x)&=0,
&
\nu(P_{\mathrm{par},\tau},x)&=0,\\
\theta(P_{\mathrm{par},\tau})&=\tau,
&
\mu_1(P_{\mathrm{par},\tau},x)
&=\ell\{\nu(P_{\mathrm{par},\tau},x)+\theta(P_{\mathrm{par},\tau})\}.
\end{aligned}
\]
The nuisance envelopes and seminorms therefore satisfy
\[
\|g(P_{\mathrm{par},\tau},\cdot)\|_\infty
=
\|\nu(P_{\mathrm{par},\tau},\cdot)\|_\infty
=
[g(P_{\mathrm{par},\tau},\cdot)]_\alpha
=
[\nu(P_{\mathrm{par},\tau},\cdot)]_\beta
=0.
\]
Together with the uniform design and homogeneous relation, these identities verify every condition of \cref{obj:def:model} whenever \(|\tau|\le1/2\).

Define the comparison amplitude and laws by
\[
\tau_n:=\frac14 n^{-1/2},
\qquad
P_{\mathrm{null}}:=P_{\mathrm{par},0},
\qquad
P_{\mathrm{shift}}:=P_{\mathrm{par},\tau_n}.
\]
Since \(n\ge1\),
\[
0\le n^{-1/2}\le1,
\qquad
0\le\tau_n\le\frac14\le\frac12.
\]
Hence, using \cref{obj:def:radius-model},
\[
P_{\mathrm{null}}\in\mathcal M_r(\alpha,\beta),
\qquad
P_{\mathrm{shift}}\in\mathcal M(\alpha,\beta),
\qquad
\theta(P_{\mathrm{null}})=0,
\qquad
\theta(P_{\mathrm{shift}})=\tau_n.
\]
Indeed, the radius condition at the null is \(0\le r\), including when \(r=0\).
% lean: CausalSmith.Stat.LogoddsLowsmoothFrontier.parametric_comparison

\item All four null cells have probability \(1/4\). The likelihood ratio of \(P_{\mathrm{par},\tau}\) relative to \(P_{\mathrm{null}}\) is therefore
\[
\Lambda_\tau(x,a,y)
:=
\frac{dP_{\mathrm{par},\tau}}{dP_{\mathrm{null}}}(x,a,y)
=
\begin{cases}
1,&a=0,\\
2\{1-\ell(\tau)\},&a=1,\ y=0,\\
2\ell(\tau),&a=1,\ y=1,
\end{cases}
\qquad \tau\in\mathbb R.
\]
This bounded density establishes absolute continuity and integrability of its squared deviation. Define its one-record chi-square divergence by
\[
\chi^2_\tau
:=
\int_\Omega(\Lambda_\tau-1)^2\,dP_{\mathrm{null}}.
\]
Summing over the four null cells gives
\[
\begin{aligned}
\chi^2_\tau
&=
\frac14\left[
\{2\ell(\tau)-1\}^2+
\{2(1-\ell(\tau))-1\}^2
\right]\\
&=2\{\ell(\tau)-1/2\}^2.
\end{aligned}
\]
For every real \(\tau\), differentiation yields
\[
0\le\ell'(\tau)
=\ell(\tau)\{1-\ell(\tau)\}
=\frac14-\{\ell(\tau)-1/2\}^2
\le\frac14.
\]
The mean-value bound, together with \(\ell(0)=1/2\), consequently gives
\[
|\ell(\tau)-1/2|\le\frac{|\tau|}{4},
\qquad
\chi^2_\tau
\le2\left(\frac{|\tau|}{4}\right)^2
=\frac{\tau^2}{8}.
\]
% lean: CausalSmith.Stat.LogoddsLowsmoothFrontier.parametricLaw_chiSq_le

\item The amplitude has the exact information budget
\[
\frac{n\tau_n^2}{8}=\frac1{128},
\qquad
n\chi^2_{\tau_n}\le\frac1{128}.
\]
Since the exponential lies above its tangent at zero,
\[
(1+\chi^2_{\tau_n})^n
\le
\{\exp(\chi^2_{\tau_n})\}^n
=
\exp(n\chi^2_{\tau_n})
\le
\exp(1/128).
\]
Comparison of the exponential and geometric series gives the required numerical bound:
\[
\exp(1/128)
=
\sum_{j=0}^{\infty}\frac{(1/128)^j}{j!}
\le
\sum_{j=0}^{\infty}(1/128)^j
=
\frac{128}{127}.
\]

For the record tuple
\[
\mathcal O=(\mathcal O_1,\ldots,\mathcal O_n)\in\Omega^n,
\]
define the product likelihood ratio and its chi-square divergence by
\[
\Lambda_{\tau,n}(\mathcal O)
:=
\prod_{j=1}^n\Lambda_\tau(\mathcal O_j),
\qquad
\chi^2_{\tau,n}
:=
\int_{\Omega^n}(\Lambda_{\tau,n}-1)^2
\,dP_{\mathrm{null}}^{\otimes n},
\qquad \tau\in\mathbb R.
\]
The product likelihood ratio is bounded and is the density of
\(P_{\mathrm{par},\tau}^{\otimes n}\) relative to
\(P_{\mathrm{null}}^{\otimes n}\). Because each likelihood ratio integrates to one, independence implies
\[
\begin{aligned}
1+\chi^2_{\tau,n}
&=
\int_{\Omega^n}\Lambda_{\tau,n}^2
\,dP_{\mathrm{null}}^{\otimes n}\\
&=
\left(\int_\Omega\Lambda_\tau^2\,dP_{\mathrm{null}}\right)^n
=
(1+\chi^2_\tau)^n.
\end{aligned}
\]
Applying this identity at \(\tau=\tau_n\) to the preceding exponential bound yields
\[
\chi^2_{\tau_n,n}
\le\frac{128}{127}-1
=\frac1{127}.
\]
% lean: CausalSmith.Stat.LogoddsLowsmoothFrontier.parametricLaw_iid_tv

\item Total variation is the supremum of the absolute differences of measurable-event probabilities. For probability laws with a density, it equals half the integral of the absolute density difference. Cauchy--Schwarz therefore gives
\[
\begin{aligned}
\operatorname{TV}
\bigl(P_{\mathrm{null}}^{\otimes n},
      P_{\mathrm{shift}}^{\otimes n}\bigr)
&=
\frac12
\int_{\Omega^n}|\Lambda_{\tau_n,n}-1|
\,dP_{\mathrm{null}}^{\otimes n}\\
&\le
\frac12\sqrt{\chi^2_{\tau_n,n}}
\le
\frac12\sqrt{\frac1{127}}
\le
\frac1{20},
\end{aligned}
\]
where the last inequality follows from \(1/127\le1/100\).

Define the full sample-and-seed experiment laws by
\[
\mathbb P_{\mathrm{null}}
:=
P_{\mathrm{null}}^{\otimes n}\otimes\lambda,
\qquad
\mathbb P_{\mathrm{shift}}
:=
P_{\mathrm{shift}}^{\otimes n}\otimes\lambda.
\]
Their likelihood ratio is still \(\Lambda_{\tau_n,n}(\mathcal O)\). Integrating over the common independent seed, whose law has mass one, gives
\[
\operatorname{TV}
(\mathbb P_{\mathrm{null}},\mathbb P_{\mathrm{shift}})
=
\operatorname{TV}
\bigl(P_{\mathrm{null}}^{\otimes n},
      P_{\mathrm{shift}}^{\otimes n}\bigr)
\le\frac1{20}.
\]
The comparison effects are ordered:
\[
\theta(P_{\mathrm{null}})
=0
\le\tau_n
=\theta(P_{\mathrm{shift}}).
\]
% lean: CausalSmith.Stat.LogoddsLowsmoothFrontier.parametricLaw_experiment_tv

\item To bound the infimum in \cref{obj:def:length-objective}, first observe that the honest-procedure class is nonempty. Define
\[
I_{\mathrm{full}}(\mathcal O,U):=[-1/2,1/2],
\qquad
(\mathcal O,U)\in\Omega^n\times[0,1].
\]
By \cref{obj:ass:effect-envelope}, for every
\(P\in\mathcal M(\alpha,\beta)\),
\[
(P^{\otimes n}\otimes\lambda)
\{\theta(P)\in I_{\mathrm{full}}(\mathcal O,U)\}=1.
\]
The constant procedure is measurable and hence belongs to the class in
\cref{obj:def:honest-procedures}.

Now fix any
\[
I\in\mathcal A_n(\alpha,\beta).
\]
For
\(\omega=(\mathcal O,U)\in\Omega^n\times[0,1]\), define the measurable coverage events
\[
\begin{aligned}
\mathcal E_{\mathrm{null}}
&:=
\{(\mathcal O,U):0\in I(\mathcal O,U)\},\\
\mathcal E_{\mathrm{shift}}
&:=
\{(\mathcal O,U):\tau_n\in I(\mathcal O,U)\}.
\end{aligned}
\]
Both comparison laws belong to the ambient model, so
\cref{obj:ass:global-coverage} gives
\[
\mathbb P_{\mathrm{null}}(\mathcal E_{\mathrm{null}})
\ge\frac9{10},
\qquad
\mathbb P_{\mathrm{shift}}(\mathcal E_{\mathrm{shift}})
\ge\frac9{10}.
\]
% lean: CausalSmith.Stat.LogoddsLowsmoothFrontier.twoPoint_lengthObjective_lower

\item The total variation bound transfers the second coverage event to the null experiment:
\[
\begin{aligned}
\mathbb P_{\mathrm{null}}(\mathcal E_{\mathrm{shift}})
&\ge
\mathbb P_{\mathrm{shift}}(\mathcal E_{\mathrm{shift}})
-
\operatorname{TV}
(\mathbb P_{\mathrm{null}},\mathbb P_{\mathrm{shift}})\\
&\ge\frac9{10}-\frac1{20}.
\end{aligned}
\]
Using the union--intersection identity and the fact that the union has probability at most one,
\[
\begin{aligned}
\mathbb P_{\mathrm{null}}
(\mathcal E_{\mathrm{null}}\cap\mathcal E_{\mathrm{shift}})
&=
\mathbb P_{\mathrm{null}}(\mathcal E_{\mathrm{null}})
+
\mathbb P_{\mathrm{null}}(\mathcal E_{\mathrm{shift}})
-
\mathbb P_{\mathrm{null}}
(\mathcal E_{\mathrm{null}}\cup\mathcal E_{\mathrm{shift}})\\
&\ge
\frac9{10}+\frac9{10}-\frac1{20}-1
=\frac34.
\end{aligned}
\]

Write \(I(\omega)=I(\mathcal O,U)\). On the intersection event, membership of both effects implies
\[
\inf I(\omega)\le0\le\tau_n\le\sup I(\omega),
\qquad
|I(\omega)|
=
\sup I(\omega)-\inf I(\omega)
\ge\tau_n.
\]
These weak endpoint inequalities apply to every choice of endpoint inclusion. Off the intersection event, interval length is nonnegative. Thus
\[
|I(\omega)|\ge
\begin{cases}
\tau_n,
&\omega\in\mathcal E_{\mathrm{null}}\cap\mathcal E_{\mathrm{shift}},\\
0,
&\omega\notin\mathcal E_{\mathrm{null}}\cap\mathcal E_{\mathrm{shift}}.
\end{cases}
\]
Admissible intervals lie in the scalar effect region, so
\[
0\le |I(\omega)|\le1.
\]
Their measurable lengths are therefore integrable. Integrating the pointwise bound yields
\[
\mathbb E_{\mathbb P_{\mathrm{null}}}|I(\mathcal O,U)|
\ge
\tau_n\,
\mathbb P_{\mathrm{null}}
(\mathcal E_{\mathrm{null}}\cap\mathcal E_{\mathrm{shift}})
\ge\frac34\tau_n.
\]
% lean: CausalSmith.Stat.LogoddsLowsmoothFrontier.twoPoint_expectedLength_lower

\item Since \(P_{\mathrm{null}}\in\mathcal M_r(\alpha,\beta)\), its expected length is bounded above by the supremum over that slice:
\[
\sup_{P\in\mathcal M_r(\alpha,\beta)}
\mathbb E_{P^{\otimes n}\otimes\lambda}|I(\mathcal O,U)|
\ge
\mathbb E_{\mathbb P_{\mathrm{null}}}|I(\mathcal O,U)|
\ge\frac34\tau_n.
\]
This supremum is finite because every expected length is at most one. The bound holds for every honest procedure, and that class is nonempty. Taking its infimum as in \cref{obj:def:length-objective} gives
\[
J_n(\alpha,\beta;r)
\ge
\frac34\{\theta(P_{\mathrm{shift}})-\theta(P_{\mathrm{null}})\}
=
\frac34\tau_n
=
\frac3{16}n^{-1/2}.
\]
% lean: CausalSmith.Stat.LogoddsLowsmoothFrontier.parametric_null_floor
\end{enumerate}
\end{proof}
